% Chapter 4 -- Evaluation (\input by main.tex): detection protocol + detection results
% + recall analysis + donor-metric validation + completion results. Merged from the
% former sec_4_1 / sec_4_results / sec_5 / sec_7_2 / sec_7_results fragments (seam
% markers below).
%
% Honesty-phrasing rules applied (personal_agenda P1):
%   - "mean IoU" always qualified as "of matched detections"
%   - recall denominator stated explicitly (>=10 points surviving preprocessing)
%   - evaluation is leave-one-sequence-out (LOSO): each sequence is scored by a
%     classifier that never saw it in training or selection; seq 08 is the reported
%     operating point (fold-08 confirms the final seq-08 numbers)
%   - measured vs inferred evidence tiers marked in prose
%   - every number traceable to a finding / output artifact
%
\chapter{Evaluation}
\label{ch:evaluation}

This chapter evaluates the full pipeline end to end, but its destination is the
contribution: the donor-frame coverage metric and what it reveals about
completion. The route there runs through the detection front-end that supplies
the metric its inputs. We first fix the detection evaluation protocol
(Section~\ref{sec:eval-protocol}) and report the detection operating point,
ablations, range, and runtime (Sections~\ref{sec:results-detection}--\ref{sec:results-runtime});
we then diagnose why detection recall is capped, because the same clustering
step that limits recall is what governs input cleanliness for completion
(Section~\ref{sec:recall-analysis}). The chapter then turns to its centre: the
donor metric is put through its validation battery
(Section~\ref{sec:donor-validation}) and used to score completion on real cars ---
surface coverage, the downstream amodal box, moving cars, and a pre-registered
held-out replication (Sections~\ref{sec:results-coverage}--\ref{sec:results-heldout}).

\section{Evaluation protocol}
\label{sec:eval-protocol}

Every quantitative claim about detection in this thesis is produced by a single
evaluation script (\texttt{src/evaluate.py}) under one fixed protocol. Because
the results, ablation, and recall-analysis sections all report numbers from this
protocol, we state it in full here and refer back to it rather than restating it
per table. The protocol was frozen before the write-up began; no configuration,
checkpoint, or metric definition described below was altered to produce a
reported number.

\subsection{Sequences, roles, and leave-one-sequence-out cross-validation}
\label{sec:eval-sequences}

We evaluate on the eleven labelled SemanticKITTI odometry sequences
(\numrange{00}{10}), whose per-point semantic and instance labels serve as
ground truth. The official test sequences \numrange{11}{21} carry \emph{no}
public labels, so no ground-truth detection metric can be computed on them; this
is a property of the dataset, not a choice, and it is why no labelled sequence
can be reserved as a conventional held-out test set.

\paragraph{Classifier evaluation uses leave-one-sequence-out cross-validation.}
Rather than report detection on the sequence the classifier was tuned against, we
evaluate the binary classifier under \emph{leave-one-sequence-out} (LOSO)
cross-validation over all eleven sequences. For each held-out sequence~$K$ a
classifier is trained on the other ten sequences and then run on the full
sequence~$K$; the held-out sequence contributes \emph{zero} clusters to that
fold's training data and is never used to select the checkpoint. Checkpoint
selection is removed entirely: every fold trains for a fixed \num{15}-epoch
budget and the last-epoch checkpoint is used (historically the best epoch was
\num{13}/\num{15}, so last~$\approx$~best), so the held-out sequence is untouched
during both training and model selection. In the per-fold table below,
therefore, the classifier scoring sequence~$K$ has never seen
sequence~$K$ during training or model selection. Because the
Stage~B clusters are stored with a sequence prefix in their filenames, the folds
are a re-partition of the existing mined data and require no re-mining.

\paragraph{Two sequences carry additional, distinct roles.}
\begin{itemize}
  \item \textbf{Sequence 08 (reported operating point; \num{4071} frames).} The
    largest labelled sequence is the operating point at which the pipeline is
    reported in detail: the stage ablations (Section~\ref{sec:results-ablations}),
    the recall root-cause analysis (Section~\ref{sec:recall-analysis}), the
    IoU-threshold sensitivity below, and the completion evaluation
    (Section~\ref{sec:donor-validation}) are all computed on the sequence-08 output
    of the \emph{final} classifier. Sequence~08 was never used to tune the
    classical stages. The final classifier was historically selected on
    sequence~08, so its sequence-08 numbers could in principle be optimistic; the
    LOSO fold-08 classifier (trained on \numrange{00}{07},\,\numrange{09}{10},
    last-epoch, never selected on 08) settles this empirically below.
  \item \textbf{Sequence 00 (completion held-out; 100 frames).} Sequence~00 is
    the held-out sequence for the \emph{completion} replication
    (Section~\ref{sec:donor-validation}), where its being unseen by the completion
    model is what matters. Its classical pipeline parameters
    (\texttt{PIPELINE\_CONFIG}) were historically hand-tuned on it, so the
    seq-00 detection fold retains classical-stage familiarity even though its
    classifier is still properly held out; we isolate that residual in the
    aggregate below (dropping fold-00 moves pooled recall by \num{0.022}).
\end{itemize}

\subsection{Ground-truth targets and the eligibility rule}
\label{sec:eval-targets}

The pipeline is a car detector, so the evaluation counts only the two
car-bearing SemanticKITTI classes as valid ground-truth objects:
\emph{car} (class~10) and \emph{moving-car} (class~252); we refer to this target
set as \emph{supported-vehicles}. A ground-truth car \emph{instance} is a
maximal set of points sharing one instance id within these classes.

An instance is counted in the evaluation only if at least \textbf{10 of its
points survive preprocessing} --- the $z$-filter, voxel downsampling, statistical
denoising, and ground removal that every scan passes through before clustering.
This is the recall denominator: \emph{recall is computed against ground-truth car
instances with $\geq 10$ points surviving preprocessing}, not against every
annotated car. The rule exists because an instance reduced to a handful of points
by occlusion or range cannot be segmented by any clustering method, so including
it would measure sensor sparsity rather than the pipeline.

This choice narrows the denominator, and honesty requires quantifying by how
much. On a stride-20 sample of sequence~08 (204 of \num{4071} frames, drawn
uniformly across the drive), of all annotated car instances
(\num{2332} pooled), \textbf{25.5\% are excluded} by the surviving-points rule;
measured instead against instances that have $\geq 10$ points in the \emph{raw}
scan (\num{2033} pooled), \textbf{14.6\% are excluded} by preprocessing.\footnote{%
These are two distinct $\geq 10$-point rules: one thresholds \emph{raw} label
points (a distinct-car count), the other thresholds points \emph{surviving
preprocessing} (the eval denominator). They are not interchangeable.}
Because an excluded instance can never be matched, the sequence-08 operating-point
recall of \num{0.730} against the eligible denominator corresponds, by inference, to a
recall of \emph{approximately} \num{0.54} against \emph{all} annotated cars
(the eligible-to-annotated survival ratio on this sample is
$\num{1737}/\num{2332}=\num{0.745}$; the ratio is a stride-20 estimate, so this
figure is inferred, not separately measured). We report the eligible-denominator
recall as the headline and disclose the all-annotated figure alongside it.

\subsection{Matching and metrics}
\label{sec:eval-matching}

Detections and ground-truth instances are compared at the \emph{point level}.
For a detection mask $D$ and an instance mask $G$ over the surviving points, the
overlap is the intersection-over-union
\[
  \mathrm{IoU}(D,G) \;=\; \frac{|D \cap G|}{|D \cup G|}.
\]
All detection--instance pairs with $\mathrm{IoU} \geq \tau$ are collected, sorted
by descending IoU, and assigned \emph{greedily and one-to-one}: the highest-IoU
pair is matched first, and once a detection or an instance is used it cannot be
matched again. This yields, per frame:
\begin{itemize}
  \item a \textbf{true positive} (TP) for each matched pair,
  \item a \textbf{false positive} (FP) for each unmatched detection,
  \item a \textbf{false negative} (FN) for each unmatched eligible instance.
\end{itemize}
Precision, recall, and $F_1$ are then computed from TP, FP, and FN
\emph{micro-averaged} over the sequence: counts are pooled across all frames and
the ratios taken once, so denser frames contribute proportionally more. We also
report the \textbf{mean IoU of matched detections} --- the mean of the IoU values
of the accepted TP pairs only. This is a quality measure of the matches that were
made; it says nothing about missed cars and must not be read as an overall
segmentation score.

Track-level filtering is enabled by default for the headline numbers: an offline
centroid tracker links detections across frames, and only tracks meeting a
minimum length and class-consistency requirement are accepted. This is a
post-hoc, offline step (it uses the whole sequence), consistent with the
offline scope of the system; its effect is isolated as an ablation row in the
ablation section (Section~\ref{sec:results-ablations}).

\subsection{Operating point}
\label{sec:eval-operating-point}

Table~\ref{tab:headline} gives the reported operating point on the full
sequence~08 under this protocol, produced by the final classifier (the one the
completion pipeline actually runs with). This is the operating point the
ablation and recall-analysis sections of this chapter decompose; the
cross-validation that confirms it follows in Section~\ref{sec:eval-loso}.

\begin{table}[H]
  \centering
  \caption[Reported detection operating point]{Reported detection operating point under the frozen protocol, using
    the final classifier. Sequence~08 is the largest labelled sequence and was
    never used to tune the classical stages; sequence~00 is a completion held-out
    reference whose classifier saw seq-00 clusters in training (so its detection
    recall reads optimistically). Mean IoU is the mean over matched detections
    only. Section~\ref{sec:eval-loso} re-scores every sequence with a classifier
    that never saw it.}
  \label{tab:headline}
  \setlength{\tabcolsep}{5pt}
  \begin{tabular}{lccccc}
    \toprule
    Sequence & Precision & Recall & $F_1$ & Mean IoU & TP / FP / FN \\
    \midrule
    08 (full, \num{4071} fr.) & 0.905 & 0.730 & 0.808 & 0.912 & \num{25478} / \num{2676} / \num{9444} \\
    00 (100 fr., reference)   & 0.967 & 0.777 & 0.862 & 0.962 & \num{1296} / \num{44} / \num{371} \\
    \bottomrule
  \end{tabular}
\end{table}

The pipeline is precision-saturated and recall-limited: false positives are flat
at roughly one per frame, while recall is set by how many eligible cars the
clustering stage recovers as single objects. The remaining sections of this
chapter report the ablations that decompose this operating point, and Section~\ref{sec:recall-analysis}
root-causes the recall shortfall.

\subsection{Leave-one-sequence-out cross-validation}
\label{sec:eval-loso}

The operating point above is reported on sequence~08, which is also the sequence
the final classifier was historically selected on. To show that this selection
does not inflate the result, and to characterise the per-sequence recall spread
that a single reporting sequence cannot, we ran the LOSO protocol of
Section~\ref{sec:eval-sequences}: eleven folds, each training a classifier on ten
sequences and evaluating the full held-out eleventh, so every fold's classifier
is blind to the sequence it scores. Table~\ref{tab:loso} reports all eleven folds,
the pooled micro-average (counts summed across folds, comparable to the single
operating-point number), and the macro mean~$\pm$~standard deviation across folds.

\begin{table}[H]
  \centering
  \caption[Leave-one-sequence-out cross-validation]{Leave-one-sequence-out cross-validation over all eleven labelled
    sequences (full sequences, $\tau=0.3$, final configuration). Each fold's
    classifier trained on the other ten sequences and was never selected on the
    held-out one, so no row's classifier has seen the sequence it scores.
    ``Pooled'' sums TP/FP/FN across folds; the excl.-seq00 pool isolates the
    residual seq-00 classical-tuning familiarity (Section~\ref{sec:eval-sequences}).}
  \label{tab:loso}
  \setlength{\tabcolsep}{3pt}
  \footnotesize
  \begin{tabular}{lccccc}
    \toprule
    Held-out seq & Precision & Recall & $F_1$ & Mean IoU & TP / FP / FN \\
    \midrule
    00 (\num{4541} fr.) & 0.913 & 0.761 & 0.830 & 0.932 & \num{42763} / \num{4079} / \num{13427} \\
    01 (\num{1101} fr.) & 0.951 & 0.336 & 0.496 & 0.953 & \num{539} / \num{28} / \num{1067} \\
    02 (\num{4661} fr.) & 0.793 & 0.700 & 0.743 & 0.903 & \num{10919} / \num{2857} / \num{4685} \\
    03 (\num{801} fr.)  & 0.733 & 0.623 & 0.673 & 0.904 & \num{1527} / \num{556} / \num{926} \\
    04 (\num{271} fr.)  & 0.797 & 0.501 & 0.615 & 0.835 & \num{529} / \num{135} / \num{527} \\
    05 (\num{2761} fr.) & 0.861 & 0.731 & 0.791 & 0.910 & \num{10671} / \num{1723} / \num{3927} \\
    06 (\num{1101} fr.) & 0.861 & 0.637 & 0.732 & 0.906 & \num{6374} / \num{1031} / \num{3637} \\
    07 (\num{1101} fr.) & 0.905 & 0.729 & 0.808 & 0.929 & \num{11131} / \num{1162} / \num{4137} \\
    08 (\num{4071} fr.) & 0.881 & 0.737 & 0.803 & 0.912 & \num{25733} / \num{3477} / \num{9189} \\
    09 (\num{1591} fr.) & 0.824 & 0.646 & 0.724 & 0.911 & \num{5141} / \num{1097} / \num{2813} \\
    10 (\num{1201} fr.) & 0.725 & 0.559 & 0.632 & 0.871 & \num{1832} / \num{694} / \num{1444} \\
    \midrule
    \textbf{Pooled (all 11)} & \textbf{0.874} & \textbf{0.719} & \textbf{0.789} & \textbf{0.919} & \num{117159} / \num{16839} / \num{45779} \\
    Pooled (excl.\ seq~00)   & 0.854 & 0.697 & 0.767 & 0.911 & \num{74396} / \num{12760} / \num{32352} \\
    Macro mean $\pm$ std     & $0.840{\pm}0.070$ & $0.633{\pm}0.121$ & $0.713{\pm}0.097$ & $0.906{\pm}0.030$ & --- \\
    \bottomrule
  \end{tabular}
\end{table}

Three things follow. First, \textbf{model selection did not inflate the reported
recall}. The fold-08 classifier, blind to sequence~08, reaches recall \num{0.737} against
the sequence-08 operating point's \num{0.730} --- if anything marginally higher,
not lower --- so the recall headline is not a product of having selected the
classifier on the reporting sequence. The model-selection optimism that does
exist is confined to precision and is small: fold-08 precision is \num{0.881}
versus the final \num{0.905} (an extra \num{801} false positives across
\num{4071} frames), exactly the two-point, precision-only exposure the stage
decomposition predicts (Section~\ref{sec:disc-eval-validity}). The two
sequence-08 rows are not two separate experiments but \emph{one classifier under
two checkpoint-selection rules}: the final model and the fold-08 model share the
same architecture, the same ten training sequences (\numrange{00}{07},
\numrange{09}{10}), and the same from-scratch \num{15}-epoch regime, and differ
only in which epoch is frozen --- the final model takes the epoch that scored
best on the sequence-08 validation set, the fold-08 model takes the last epoch and
never inspects sequence~08 at all. The \num{0.730}~$\rightarrow$~\num{0.737}
(recall) and \num{0.905}~$\rightarrow$~\num{0.881} (precision) deltas therefore
isolate the effect of that selection: selecting on sequence~08 bought about two
precision points and no recall. The isolation is not perfect --- because
sequence~08 cannot serve as the fold's loss-monitoring split, the fold-08 model
trains on \SI{95}{\percent} of the final model's clusters (a \num{5}-percent
monitoring slice is set aside) --- but that confound removes training data from
the fold-08 model, which would depress its recall, not inflate it, so it only
strengthens the ``recall was not inflated'' reading. Second, the
\textbf{pooled cross-validated operating point} over all eleven sequences ---
precision \num{0.874}, recall \num{0.719}, $F_1$~\num{0.789}, mean IoU
\num{0.919} --- sits within roughly one point of the sequence-08 operating point
on every metric, so the reported numbers are representative of the labelled set
and not specific to sequence~08. Third, the \textbf{per-sequence spread} is now
explicit: recall ranges from \num{0.336} on sequence~01 (a sparse highway drive)
to \num{0.761} on sequence~00 (dense urban), with a macro mean of
\num{0.633}~$\pm$~\num{0.121}; precision is far steadier
(\num{0.840}~$\pm$~\num{0.070}). The gap between the pooled recall (\num{0.719})
and the macro mean (\num{0.633}) is driven by the smallest, sparsest sequences
(seq~01 and seq~04, the latter only \num{271} frames), which the count-weighted
pool down-weights and the unweighted mean does not; both figures are reported so
neither hides the spread. Finally, the residual seq-00 classical-tuning exposure
is bounded: dropping the seq-00 fold lowers pooled recall only from \num{0.719}
to \num{0.697} (a \num{0.022} delta), so the historical seq-00 pipeline tuning
does not drive the aggregate.

\subsection{Choice of IoU threshold and its sensitivity}
\label{sec:eval-iou-threshold}

All headline numbers use $\tau = 0.3$. A point-level threshold of 0.3 is
permissive relative to the bounding-box 0.5 common in detection benchmarks, so
it is fair to ask whether the results are an artifact of a lenient gate. They
are not, and the reason is visible in Table~\ref{tab:headline}: the mean IoU of
matched detections is \num{0.912} on sequence~08, far above either threshold, so
the matches are not marginal overlaps sitting just past $\tau$.

To make this concrete we re-ran the full sequence~08 evaluation at
$\tau \in \{0.25, 0.30, 0.50\}$ with everything else frozen
(Table~\ref{tab:iou-sensitivity}).

\begin{table}[H]
  \centering
  \caption[IoU-threshold sensitivity on seq~08]{IoU-threshold sensitivity on the full sequence~08 (frozen
    configuration). The $\tau = 0.30$ row is the headline. $F_1$ moves by at most
    \num{2.6} points across the tested range (\num{0.811} at $\tau=0.25$ to
    \num{0.785} at $\tau=0.50$); mean IoU of matched detections rises at the
    stricter gate because it keeps only the better-overlapping matches.}
  \label{tab:iou-sensitivity}
  \begin{tabular}{cccccc}
    \toprule
    $\tau$ & Precision & Recall & $F_1$ & Mean IoU & TP / FP / FN \\
    \midrule
    0.25 & 0.908 & 0.732 & 0.811 & 0.910 & \num{25576} / \num{2578} / \num{9346} \\
    0.30 & 0.905 & 0.730 & 0.808 & 0.912 & \num{25478} / \num{2676} / \num{9444} \\
    0.50 & 0.879 & 0.709 & 0.785 & 0.927 & \num{24744} / \num{3410} / \num{10178} \\
    \bottomrule
  \end{tabular}
\end{table}

$F_1$ falls only \num{2.3} points from the headline $\tau = 0.30$ to the strict
$\tau = 0.50$ (\num{0.808} to \num{0.785}), and \num{2.6} points across the full
tested range from the permissive $\tau = 0.25$ (\num{0.811}), so the reported
quality is stable under the choice of threshold and does not depend on the lenient
value. We had
pre-registered that a threshold-sensitive result would be \emph{reported as a
finding rather than tuned away}; it was not sensitive.


%% ===== merged from sec_4_results.tex ===== %%

% Detection Results, Ablations, Range, Runtime (Sections 4.2--4.5; folds under Chapter 4).
% Opens with \section{Detection Evaluation} carrying its own distinct label
% (\label{ch:detection-eval-results}), which does not clash with the chapter's
% \label{ch:evaluation}.
%
% Honesty-phrasing rules applied (personal_agenda P1):
%   - "mean IoU" always qualified as "of matched detections"
%   - recall denominator restated / cross-referenced (>=10 surviving points)
%   - seq 08 = classifier model-selection set, not a held-out test set
%   - negative results reported AS FINDINGS, not apologised for
%     (C5 synthetic-pretraining redundant; clustering "adopt nothing")
%   - the distance table is per-frame / track-filter-OFF: stated as a deviation,
%     read as a trend, NOT a restatement of the 0.730 headline
%   - every number traces to a finding / output artifact (cited Finding~\#NN)
%
\section{Detection results}
\label{sec:results-detection}

The headline operating point is the sequence~08 row of Table~\ref{tab:headline}
(Section~\ref{sec:eval-protocol}): precision \num{0.905}, recall \num{0.730}, $F_1$ \num{0.808}, and
a mean IoU of matched detections of \num{0.912}, over
\num{25478}~/~\num{2676}~/~\num{9444} TP~/~FP~/~FN. Two properties of this
operating point organise the rest of the chapter, and both are visible directly
in the per-frame record rather than only in the pooled totals.

First, the pipeline is \emph{precision-saturated}. Across the full sequence the
false-positive count is low and stable --- about \num{0.7} per frame
(\num{2676} over \num{4071} frames), independent of how many cars are in view.
Figure~\ref{fig:bev-detections} shows this spatially --- detections track the
true parked and moving cars along the street, with few spurious boxes --- and
Figure~\ref{fig:timeseries} shows it over time: the FP trace is low and
stationary while the number of ground-truth cars per frame swings widely.

Second, the pipeline is \emph{recall-limited}, and the limit is structural. The
recall shortfall traces to a measured clustering failure: HDBSCAN splits large,
close cars into fragments that no longer match a single instance, and the split rate is highest exactly where cars are densest and
nearest (\SI{66}{\percent} at \SIrange{0}{10}{\metre}).
Section~\ref{sec:results-distance} exhibits the resulting near-range recall dip
directly, and the time series (Figure~\ref{fig:timeseries}) shows per-frame
recall tending to fall in the denser frames while false positives stay flat. We
root-cause the splitting mechanism in Section~\ref{sec:recall-analysis}; here we establish only that
recall, not precision, is the dominant error mode, and that it does not come from
the classifier misfiring (false positives stay flat regardless).

\begin{figure}[htbp]
  \centering
  \includegraphics[width=\linewidth]{seq08_bev_detections.png}
  \caption[Bird's-eye view of detections on seq~08]{Bird's-eye view of detections on a representative sequence~08 frame.
    Detected car clusters (coloured) are overlaid on the ground-removed point
    cloud; ego is at the origin. Detections concentrate on true vehicles, with
    few spurious clusters --- the precision-saturated behaviour reported in
    Table~\ref{tab:headline}.}
  \label{fig:bev-detections}
\end{figure}

\begin{figure}[htbp]
  \centering
  \includegraphics[width=\linewidth]{seq08_timeseries.png}
  \caption[Per-frame detection behaviour on seq~08]{Per-frame detection behaviour over sequence~08. False positives (lower
    trace) stay low and flat --- under one per frame --- regardless of scene
    content, while per-frame recall tends to fall in frames with more
    ground-truth cars. The structural recall limit is root-caused in Section~\ref{sec:recall-analysis}
    (HDBSCAN splitting).}
  \label{fig:timeseries}
\end{figure}

\section{Ablations}
\label{sec:results-ablations}

This section decomposes the headline operating point into the contribution of
each learned or post-hoc stage. Every ablation row is produced under the frozen
protocol of Section~\ref{sec:eval-protocol} unless a deviation is stated explicitly.

\subsection{Stage decomposition: classifier and track-filter contributions}
\label{sec:ablation-stage}

Table~\ref{tab:ablation-stage} isolates the two filtering stages that sit on top
of geometric clustering. The geometric-only row uses clustering plus the
geometric size/shape filters with \emph{no} learned classifier and \emph{no}
track filter (like-for-like under the promoted
configuration); the full row is the headline.

\begin{table}[H]
  \centering
  \caption[Stage ablation on seq~08]{Stage ablation on the full sequence~08 (frozen configuration). Rows
    add one filtering stage at a time on top of geometric clustering. The
    per-frame classifier is the precision mechanism (it removes $\sim$\num{149000}
    junk clusters) but is conservative on borderline cars (recall drops); the
    offline track filter then recovers recall by temporal voting while raising
    precision further. Mean IoU is over matched detections only. Geometric-only
    and the middle row are per-frame (no track filter); the full row uses the
    offline track filter.}
  \label{tab:ablation-stage}
  \setlength{\tabcolsep}{5pt}
  \begin{tabular}{lccccc}
    \toprule
    Configuration & Precision & Recall & $F_1$ & Mean IoU & TP / FP / FN \\
    \midrule
    Geometric only                & 0.149 & 0.775 & 0.250 & 0.907 & \num{27051} / \num{154612} / \num{7871} \\
    \;+\,classifier              & 0.820 & 0.677 & 0.741 & 0.921 & \num{23629} / \num{5204} / \num{11293} \\
    \;+\,classifier, track filter & 0.905 & 0.730 & 0.808 & 0.912 & \num{25478} / \num{2676} / \num{9444} \\
    \bottomrule
  \end{tabular}
\end{table}

The decomposition tells a two-step story. The \textbf{classifier
is the precision mechanism}: geometric clustering under the promoted coarse-voxel
setting emits many mostly-junk clusters (precision \num{0.149}), and the per-frame
classifier removes about \num{149000} of them across the sequence
(FP \num{154612}~$\rightarrow$~\num{5204}), lifting precision from \num{0.149} to
\num{0.820}. It pays for this with recall --- being conservative, it discards
borderline true clusters too, so recall falls from \num{0.775} to \num{0.677}.
The \textbf{track filter then recovers that recall while improving precision
again}: linking detections across frames and accepting only consistent tracks
lifts precision to \num{0.905} \emph{and} recall to \num{0.730} (TP
\num{23629}~$\rightarrow$~\num{25478}), because a car seen across many frames is
kept even where a single frame classified it as unknown, while transient
false positives fail the track-length test. The net of both stages is still a
recall trade --- geometric-only recall (\num{0.775}) exceeds the full-pipeline
recall (\num{0.730}), roughly \num{1600} true positives forgone for the
$\sim$\num{152000}-false-positive removal --- but that trade runs through a
recall dip and a partial recovery, not a simple monotone loss. We report it
rather than hiding it: the system is deliberately tuned for precision, and the
recall it forgoes here is small next to the recall it never had (Section~\ref{sec:recall-analysis}). The
pre-promotion geometric-only figure ($F_1$~\num{0.305}) is
superseded by this like-for-like \num{0.250} and kept only as a historical
record.

\subsection{Redundancy of synthetic pretraining (Stage~A ablation)}
\label{sec:ablation-stagea}

The classifier ships as a binary car/not-car network trained \emph{from scratch}
on mined real clusters (Stage~B). A natural alternative is a two-stage
recipe in which a synthetic-ShapeNet Stage~A pretraining seeds Stage~B; we tested
whether that pretraining is necessary and found it is not --- a negative result we
report as a finding, because it changes the recommended recipe.

At the classifier level, on the sequence~08 validation clusters, the from-scratch
network matches the pretrained one: macro $F_1$ \num{0.9285} from scratch versus
\num{0.9225} pretrained (the two are within run-to-run noise, with
scratch marginally ahead). At the pipeline level a small precision edge for the
pretrained checkpoint appeared on a 100-frame sample (FP \num{27} versus
\num{62}), but it was uncontrolled (different training seed and code) and
\textbf{did not survive at full scale}: on the full sequence~08 the two
checkpoints are identical in $F_1$ to three decimals. We therefore
drop Stage~A from the production pipeline and keep it as ablation material.

The reason the synthetic prior is redundant --- and cannot simply be added for
safety --- is a total, symmetric domain gap between synthetic and real point
clusters. Table~\ref{tab:cross-domain} gives the cross-domain
transfer matrix: a model trained on one domain and evaluated on the other scores
a car-class $F_1$ of \num{0.000} in \emph{both} off-diagonal directions, and
fine-tuning on real data erases whatever the synthetic pretraining had learned.
Given \num{420000} real training clusters, the real data does essentially all the
work; the synthetic prior neither helps nor, once fine-tuned, survives.

\begin{table}[H]
  \centering
  \caption[Cross-domain classifier transfer]{Cross-domain classifier transfer (car-class $F_1$). Each
    model is trained on one domain and evaluated on the other's validation set;
    diagonal cells (italic) are in-domain, for scale. Both off-diagonal cells
    collapse to \num{0.000} --- against strong in-domain performance
    (\num{0.999}, \num{0.885}) --- so synthetic and real clusters are mutually
    out-of-distribution and synthetic pretraining transfers nothing to real
    inference.}
  \label{tab:cross-domain}
  \begin{tabular}{lcc}
    \toprule
    & \multicolumn{2}{c}{Evaluated on} \\
    \cmidrule(lr){2-3}
    Trained on & Synthetic & Real \\
    \midrule
    Synthetic (Stage~A) & \emph{0.999} & 0.000 \\
    Real (Stage~B)      & 0.000        & \emph{0.885} \\
    \bottomrule
  \end{tabular}
\end{table}

\subsection{Clustering method: cost and benefit of HDBSCAN}
\label{sec:ablation-clustering}

The clustering stage is HDBSCAN. Because it is by far the most expensive stage
(Section~\ref{sec:results-runtime}), we benchmarked it against two cheaper
density/connectivity clusterers --- DBSCAN and Euclidean-tolerance clustering ---
at the same voxel resolution and under an identical pipeline
(Table~\ref{tab:clustering}). \textbf{Deviation:} this benchmark is
per-frame with the track filter \emph{off}, because it strides across the
sequence and the tracker requires contiguous frames; the numbers are therefore
comparable to each other but not to the track-filtered headline.

\begin{table}[H]
  \centering
  \caption[Clustering-method benchmark]{Clustering-method benchmark at fixed voxel resolution
    ($\SI{0.10}{\metre}$), identical pipeline, per-frame / track-filter-off. HDBSCAN wins $F_1$ on both sequences, entirely through
    recall; precision is essentially method-independent. The alternatives are
    5--7$\times$ faster. Mean IoU is over matched detections; the timing column is
    the mean per-frame clustering time.}
  \label{tab:clustering}
  \begin{tabular}{llccccc}
    \toprule
    Seq & Method & Precision & Recall & $F_1$ & Mean IoU & Clustering (mean ms) \\
    \midrule
    \multirow{3}{*}{00}
      & HDBSCAN   & 0.962 & \textbf{0.731} & \textbf{0.831} & 0.968 & 391 \\
      & DBSCAN    & 0.962 & 0.637 & 0.767 & 0.964 & 59 \\
      & Euclidean & 0.967 & 0.636 & 0.768 & 0.982 & 78 \\
    \midrule
    \multirow{3}{*}{08}
      & HDBSCAN   & 0.813 & \textbf{0.671} & \textbf{0.735} & 0.923 & 451 \\
      & DBSCAN    & 0.848 & 0.565 & 0.678 & 0.927 & 69 \\
      & Euclidean & 0.838 & 0.576 & 0.683 & 0.950 & 93 \\
    \bottomrule
  \end{tabular}
\end{table}

The gap is entirely a recall gap --- precision is within \num{0.035} across all
methods on either sequence, while HDBSCAN's recall exceeds the alternatives by
\numrange{9}{11} points. An $\varepsilon$-sweep confirms no fixed radius closes
it: the best fixed-radius setting (Euclidean at $\varepsilon=\SI{0.4}{\metre}$)
reaches $F_1$ \num{0.800} on sequence~00, still short of HDBSCAN's \num{0.831}.
HDBSCAN's variable-density criterion recovers cars that a single radius cannot,
which is exactly the recall margin that matters here. We therefore
\textbf{adopt nothing} from the benchmark and keep HDBSCAN in production, paying
its runtime cost knowingly; DBSCAN and Euclidean remain available as
command-line options for reproduction.

\section{Detection range}
\label{sec:results-distance}

Recall is not uniform in range, and the way it varies is not the naive
``far~=~worse'' monotone one might expect. Table~\ref{tab:distance} stratifies
detection by the range of each cluster's centroid from the ego sensor, on a
stride-20 sample of sequence~08 under the frozen configuration.
\textbf{Deviation:} like the clustering benchmark, this table is
\emph{per-frame with the track filter off} (striding breaks the tracker), so its
pooled recall (\num{0.671}) is a per-frame figure and does \emph{not} restate the
track-filtered headline \num{0.730}; read the table as a range \emph{trend}, not
as a re-derivation of the headline.

\begin{table}[H]
  \centering
  \caption[Detection stratified by range-to-ego on seq~08]{Detection stratified by range-to-ego on sequence~08 (stride-20 sample,
    frozen configuration, per-frame / track-filter-off). Recall is
    \emph{non-monotonic}: it dips in the nearest bin, peaks at
    \SIrange{10}{20}{\metre}, then falls with range. The nearest-bin dip is
    over-segmentation, the far decline is sparsity. Mean IoU is
    over matched detections. Summing the two near bins recovers the
    \SIrange{0}{20}{\metre} figure ($R=\num{0.829}$) comparable to the earlier
    range analysis.}
  \label{tab:distance}
  \begin{tabular}{lccccc}
    \toprule
    Range (m) & Precision & Recall & $F_1$ & Mean IoU (matched) & TP / FP / FN \\
    \midrule
    0--10   & 0.896 & 0.787 & 0.838 & 0.902 & \num{311} / \num{36}  / \num{84}  \\
    10--20  & 0.855 & \textbf{0.862} & 0.859 & 0.926 & \num{432} / \num{73}  / \num{69}  \\
    20--40  & 0.762 & 0.549 & 0.638 & 0.934 & \num{400} / \num{125} / \num{329} \\
    40+     & 0.393 & 0.196 & 0.262 & 0.943 & \num{22}  / \num{34}  / \num{90}  \\
    \midrule
    Pooled  & 0.813 & 0.671 & 0.735 & 0.923 & \num{1165} / \num{268} / \num{572} \\
    \bottomrule
  \end{tabular}
\end{table}

The counter-intuitive result is the \textbf{nearest-bin recall dip}: recall at
\SIrange{0}{10}{\metre} (\num{0.787}) is \emph{lower} than at
\SIrange{10}{20}{\metre} (\num{0.862}), even though near cars carry far more
points. This is not sparsity --- it is over-segmentation. The split-rate analysis shows
that HDBSCAN splits the majority of the very-near cars (a split rate of
\SI{66}{\percent} at \SIrange{0}{10}{\metre}, falling to \SI{4}{\percent} at
\SIrange{30}{50}{\metre}): a large, close car dissolves into several clusters,
none of which matches the whole instance, so it is counted as a miss. The nearest bin
also has the \emph{lowest} mean IoU of matched detections (\num{0.902}),
consistent with fragments eroding the overlap even on the cars that do match. We
develop this splitting mechanism, and why post-hoc fragment-merging fails to
recover it, in Section~\ref{sec:recall-analysis}.

Beyond the peak the decline is ordinary sparsity: recall falls from \num{0.862}
at \SIrange{10}{20}{\metre} to \num{0.549} at \SIrange{20}{40}{\metre} and
\num{0.196} past \SI{40}{\metre}, as cars thin to too few surviving points to
cluster or to clear the \num{10}-point eligibility rule (Section~\ref{sec:eval-protocol}).
Mean IoU of matched detections \emph{rises} with range
(\num{0.902}~$\rightarrow$~\num{0.943}): the far cars that are detected at all are
compact, single-cluster instances, so their overlap is clean even as their number
collapses. The earlier precision-only, seq-00, pre-promotion range analysis is superseded by this on-frozen-config sequence-08 table and is
retained only as a historical note.

\section{Runtime}
\label{sec:results-runtime}

Table~\ref{tab:runtime} breaks the per-frame cost down by stage, measured over
the full sequence~08 under the final (promoted) configuration
(\num{4068} timed frames, RTX~3070~Ti / Ryzen~7~7800X3D;
per-stage figures from \texttt{timing\_seq08\_full\_n4068\_combined.json}). The breakdown shows where
the time goes: the two classical geometric stages --- ground-plane RANSAC and,
above all, HDBSCAN clustering --- dominate, together taking roughly
\SI{68}{\percent} of the frame; adding preprocessing and the geometric filter,
the classical CPU stages account for about \SI{94}{\percent}, leaving the single
learned component --- the classifier --- at only $\sim$\SI{5.4}{\percent} (the
centroid/Hungarian tracker is a classical assignment step, not learned). The
classifier is \SI{33.7}{\milli\second} for the $\sim$\num{45} clusters in a frame
($\sim$\SI{0.75}{\milli\second} each) and the tracker is negligible.

\begin{table}[H]
  \centering
  \caption[Per-frame runtime by stage on seq~08]{Per-frame runtime by stage on the full sequence~08 under the final
    (promoted) configuration (\num{4068} timed frames; mean ms).
    Clustering alone is nearly \SI{60}{\percent} of the frame. The total is
    $\approx\SI{623.5}{\milli\second}$ ($\approx$\num{1.6}~fps) on this full run;
    a stride-20 sample of the same configuration estimates
    $\approx\SI{650.6}{\milli\second}$ (the two agree to within the sampling
    variation of the expensive stages). The pre-optimisation configuration
    totalled \SI{921.3}{\milli\second}. Percentages are of the
    \SI{623.5}{\milli\second} total.}
  \label{tab:runtime}
  \begin{tabular}{lcc}
    \toprule
    Stage & Mean (ms) & Share \\
    \midrule
    Data load                       & 0.7   & 0.1\% \\
    Preprocessing (z-filter, denoise, voxel) & 125.9 & 20.2\% \\
    Ground removal (RANSAC)         & 56.2  & 9.0\% \\
    Clustering (HDBSCAN)            & 370.4 & 59.4\% \\
    Geometric filter                & 36.2  & 5.8\% \\
    Classification                  & 33.7  & 5.4\% \\
    Tracking                        & 0.3   & 0.1\% \\
    \midrule
    Total                           & 623.5 & 100\% ($\approx$1.6 fps) \\
    \bottomrule
  \end{tabular}
\end{table}

The promoted configuration is itself the product of a runtime optimisation that
targeted the dominant classical stages directly: reducing the
RANSAC iteration count and coarsening the clustering voxel to \SI{0.10}{\metre}
cut the per-frame total by roughly a third, from a pre-optimisation
\SI{921.3}{\milli\second} to \SI{623.5}{\milli\second} on the full sequence
(\SI{650.6}{\milli\second} on the stride-20 sample), and the coarser voxel also
\emph{improved} detection --- it is the promoted configuration used throughout
this chapter. A sub-\SI{400}{\milli\second} target was shown to be unreachable
through the authorised optimisation tiers --- the post-ground object cloud is
intrinsically sparse, so coarser voxelisation compresses it little and the Python
HDBSCAN retains a floor of $\sim$\SI{185}{\milli\second} --- and the achieved
operating point was accepted, real-time operation being out of scope (user
decision, \num{2026}-07-23). Point-completion, when it runs, adds
$\sim$\SI{13}{\milli\second} per completed car on top of detection and is not part
of the detection budget above.


%% ===== merged from sec_5_recall_bottleneck.tex ===== %%

% The Recall Bottleneck (folds under Chapter 4 as a section, not a chapter of its own).
% Root-causes the recall shortfall established earlier in the chapter and reports the
% failed repair strategies as a negative-results contribution. Sub-tables are
% referenced by section + Finding number rather than hard-coded table number, to stay
% robust to renumbering.
%
% Honesty-phrasing rules applied (personal_agenda P1):
%   - recall denominator restated / cross-referenced (>=10 surviving points, Sec 4.1)
%   - "mean IoU" always qualified as "of matched detections"
%   - negative results (#21/#24, temporal aggregation) reported AS FINDINGS, not
%     apologised for
%   - the retracted "~0.74 hard limit" is corrected explicitly (#24 Correction / #34)
%   - two measurement eras kept separate: the pre-promotion seq-00 repair search
%     (baseline F1 0.844; #21/#24/temporal) vs the promoted-config seq-08 fix (#34)
%   - every number traces to a finding / session log artifact (cited Finding~\#NN /
%     session_history.md)
%
\section{Recall analysis: the bridge to completion}
\label{sec:recall-analysis}

The detection operating point above established two facts: the pipeline is
precision-saturated (false positives are low and flat, about \num{0.7} per frame),
and it is recall-limited, missing roughly one car in four
($R=\num{0.730}$ on sequence~08). This section explains \emph{why} the recall is
limited, and it is the bridge to the completion evaluation that follows: the same
clustering step that caps recall is the one that decides whether a detected car
reaches completion as a clean single-object cluster. We show that the shortfall is
a clustering failure rather than a classification failure
(Section~\ref{sec:recall-diagnosis}), root-cause it to HDBSCAN splitting large,
close cars into fragments (Section~\ref{sec:split-analysis}), report a search of
repair strategies that all fail to generalise (Section~\ref{sec:repairs}), and
describe the one intervention that did help --- coarsening the clustering
resolution --- which forces a correction of an earlier, overstated ``hard
ceiling'' claim (Section~\ref{sec:resolution-fix}). The negative results are
reported as findings, not apologised for: mapping which repairs fail, and why, is
itself a finding worth reporting, and it is what justifies the residual limitation carried to
Chapter~\ref{ch:conclusion}. Unless stated otherwise, the recall denominator
throughout is the one defined in Section~\ref{sec:eval-protocol} --- ground-truth
car instances with at least \num{10} points surviving preprocessing.

\subsection{The shortfall as a clustering failure}
\label{sec:recall-diagnosis}

The recall ceiling is upstream of the learned classifier. The stage decomposition
of Section~\ref{sec:results-ablations} makes this direct: geometric clustering with the
size/shape filters but \emph{no} learned classifier and \emph{no} track filter
already reaches recall \num{0.775}, which is \emph{higher} than the full pipeline's
\num{0.730}. Every learned stage that follows can only hold or reduce that recall:
the classifier discards clusters, and the centroid tracker only re-labels and
filters clusters that geometric clustering has already produced --- it emits a
detection in a frame only where a cluster exists in that frame, and never
interpolates a detection across a gap (\texttt{src/tracker.py};
\texttt{src/evaluate.py}), so neither stage can invent a detection where clustering
found nothing. Across the full sequence the two stages together trade roughly
\num{1600} true positives (\num{27051} geometric-only versus \num{25478} full) for the removal of $\sim$\num{152000} false positives; this is a
distinct trade from the $+\num{1655}$ true positives that the resolution change of
Section~\ref{sec:resolution-fix} later recovers. Whatever caps recall is therefore
already present in the clusters that geometric clustering produces, before the
classifier sees them.

The geometric filters are not the cap either. Table-level ablation of the
size/shape gates (sequence~00, \num{100} frames) shows that of the
\num{16454} HDBSCAN clusters in the sample, only \num{2651} (\SI{16.1}{\percent})
pass the filters, and \num{1827} of the rejected clusters do overlap a ground-truth
car --- dominated by \texttt{min\_volume} (\SI{67.9}{\percent} of the GT-matching
rejections) and \texttt{min\_points} (\SI{26.4}{\percent}); sequence~08 shows the
same pattern (\SI{66.5}{\percent} / \SI{28.0}{\percent}). The dominant gates are
\texttt{min\_volume} and \texttt{min\_points} --- the two that reject \emph{small}
clusters. This is consistent with the rejected clusters being the smaller
sub-fragments of cars that HDBSCAN has already split (the dominant fragment passing
while the splinters fail the size gates), which the split analysis of
Section~\ref{sec:split-analysis} then confirms is the prevailing failure mode. We
state this as an inference rather than a direct measurement: the ablation records
only that a rejected cluster overlaps \emph{some} ground-truth car
(\texttt{src/analyze\_clustering.py}), not whether that same car also retains a
surviving cluster, so it establishes that the rejections are dominated by small
fragments, not that every one belongs to an otherwise-detected car. Either way the
diagnosis points one stage further upstream, at the clustering itself.

\subsection{Root cause: HDBSCAN fragmentation of large, close cars}
\label{sec:split-analysis}

To quantify how clustering caps recall we measured, for every eligible ground-truth
car, how its points distribute across HDBSCAN clusters. Each car is
assigned to exactly one of four states that partition the population: \emph{ok} (one
dominant cluster, under \SI{30}{\percent} noise), \emph{split} (points spread over
two or more clusters), \emph{noisy} (one cluster but at or above \SI{30}{\percent}
noise), or \emph{all-noise}. A separate, cross-cutting flag marks \emph{merged} cars
(a car whose dominant cluster is shared with another car); a merged car still has a
dominant cluster, so it is \emph{not} a fifth partition class and is reported below
as a memo count rather than a table row that would be summed.
Table~\ref{tab:split} gives the result on both sequences.

\begin{table}[H]
  \centering
  \caption[Ground-truth cars across HDBSCAN clusters]{Distribution of ground-truth cars across HDBSCAN clusters
    (pre-promotion configuration; sequence~00, \num{100} frames;
    sequence~08, held-out sample). The four states partition the population and sum
    to the denominator; \emph{split} is the dominant failure mode and \emph{merging}
    (the cross-cutting memo row, not summed) is negligible. Split cars are uneven ---
    the largest fragment holds a median \num{0.92} of the points on sequence~00 and
    \num{0.77} on sequence~08 --- but a large point-fraction does not guarantee a
    match (see text): recall shows most split cars are lost.}
  \label{tab:split}
  \begin{tabular}{lcccc}
    \toprule
    & \multicolumn{2}{c}{Seq 00 (\num{1631} cars)} & \multicolumn{2}{c}{Seq 08 (\num{811} cars)} \\
    \cmidrule(lr){2-3}\cmidrule(lr){4-5}
    Status & Count & \% & Count & \% \\
    \midrule
    ok (single dominant cluster) & 1117 & 68.5 & 509 & 62.8 \\
    split (2+ clusters)          & 514  & 31.5 & 301 & 37.1 \\
    noisy / all-noise            & 0    & 0.0  & 1   & 0.1  \\
    \midrule
    partition total              & 1631 & 100  & 811 & 100  \\
    \addlinespace
    \textit{merged} (cross-cutting) & 8 & 0.5 & 0 & 0.0 \\
    \bottomrule
  \end{tabular}
\end{table}

Three things follow. First, \textbf{splitting, not merging, is the failure mode}:
\SIrange{31}{37}{\percent} of cars split, while merging is at most
\SI{0.5}{\percent}. The pipeline almost never fuses two cars into one detection; it
repeatedly breaks one car into several. Second, \textbf{the split is worst where
cars are largest and nearest}: the split rate falls monotonically with range ---
\SI{66}{\percent} at \SIrange{0}{10}{\metre}, \SI{40}{\percent} at
\SIrange{10}{20}{\metre}, \SI{19}{\percent} at \SIrange{20}{30}{\metre}, and
\SI{4}{\percent} at \SIrange{30}{50}{\metre}. Large, close cars return many points, and
HDBSCAN finds internal density variations within that dense return --- the roof line
against the body, the near side against the far side --- and cuts the car along
them. This is the mechanism behind the counter-intuitive near-range recall dip
reported in Section~\ref{sec:results-distance}: the nearest bin has \emph{lower} recall than
the \SIrange{10}{20}{\metre} bin despite carrying far more points, because a close
car dissolves into fragments that no longer match the whole instance. Third,
\textbf{the cleanly-clustered fraction tracks recall, and splitting usually loses
the car}. The fraction of cars that form one dominant cluster
(\SI{68.5}{\percent} on sequence~00, \SI{62.8}{\percent} on sequence~08) sits just
\emph{below} the recall the pipeline achieves on the same configuration
(\num{0.739} and \num{0.699} respectively). Two things follow. The pipeline recovers
essentially every cleanly-clustered car, so that fraction is close to a floor on
recall --- not a ceiling; recall exceeds it only by the few split cars whose dominant
fragment happens to survive matching and filtering. And because recall
(\num{0.699}) barely exceeds the clean fraction (\SI{62.8}{\percent}) while
\SI{37.1}{\percent} of cars split, the great majority of split cars are \emph{not}
recovered: a large dominant point-fraction (median \num{0.77}--\num{0.92}) does not
translate into a match, because matching depends on spatial overlap and on the
fragment clearing the size filters, not on point count. Splitting is therefore the
mechanism by which cars are lost, which is exactly why it is the recall bottleneck.
This also predicts the resolution fix of Section~\ref{sec:resolution-fix}: coarsening
the voxel converts split cars into clean ones, which should lift recall by shrinking
the reliance on recovering splits --- and it does, narrowing the sequence-08 gap
between the clean fraction and recall from \num{7.1} points (\SI{62.8}{\percent}
versus \num{0.699}) to \num{3.1} points (\SI{69.9}{\percent} versus \num{0.730}).

Figure~\ref{fig:failure-zooms} shows the mechanism directly on three sequence~08
frames. Missed cars (orange) appear as broken rows of points close to the ego
vehicle, each row too small on its own to survive the size filters, while cleaner
mid- and far-range cars (green) are detected as single clusters.

\begin{figure}[htbp]
  \centering
  \includegraphics[width=\linewidth]{seq08_failure_zooms.png}
  \caption[Detection failure close-ups on seq~08]{Failure close-ups on sequence~08 (true positives green, false positives
    red, missed ground-truth cars orange; dashed outlines are ground-truth car
    footprints). Near-ego missed cars are fragmented into several small point groups
    rather than one cluster, matching the split mechanism of
    Table~\ref{tab:split}. Although the panel title mentions both merge and split,
    merging is measured at $\leq\!\SI{0.5}{\percent}$ (Table~\ref{tab:split}), so the
    misses shown are fragmentation, not fusion of neighbouring cars.}
  \label{fig:failure-zooms}
\end{figure}

\subsection{Limits of the repair strategies}
\label{sec:repairs}

Given that splitting caps recall, the obvious responses fall into two families:
change how clusters \emph{form} (a larger minimum cluster size, a different
clusterer, distance-adaptive parameters, or pre-clustering point accumulation) or
\emph{reassemble} fragments after HDBSCAN has split them. We tried six in total;
Table~\ref{tab:strategies} collects the four with a reproducible quantitative
outcome, and two further strategies (temporal accumulation and threshold lowering)
are discussed in prose below. All were evaluated against the same pre-promotion
sequence-00 baseline (full pipeline, classifier and track filter on;
$F_1=\num{0.844}$), and the load-bearing ones were re-checked on the held-out
sequence~08.

\begin{table}[H]
  \centering
  \caption[Recall-repair strategies]{Recall-repair strategies against the pre-promotion baseline (full
    pipeline, sequence~00, \num{100} frames; held-out sequence-08 $F_1$ where
    measured). Every strategy either loses $F_1$ outright or fails to hold its
    sequence-00 gain on the held-out sequence. This is the \emph{pre-promotion}
    repair search --- it predates the coarse-voxel configuration of
    Section~\ref{sec:resolution-fix}, which is why the baseline here is
    $F_1=\num{0.844}$, not the promoted headline. Sources in the ``Src'' column are
    finding numbers. Two further strategies --- temporal point accumulation and
    lowering the size thresholds --- are discussed in the text but omitted here:
    temporal accumulation is single-sourced to a reverted feature in the session
    log (not reproducible under the research freeze), and threshold lowering
    produced no change in true positives; both belong in prose rather than the main
    results table.}
  \label{tab:strategies}
  \setlength{\tabcolsep}{4pt}
  \begin{tabular}{lcccc}
    \toprule
    Strategy & Seq 00 $P$ & Seq 00 $R$ & Seq 00 $F_1$ & Seq 08 $F_1$ \\
    \midrule
    Baseline (HDBSCAN, MCS \num{10}) & 0.984 & 0.739 & \textbf{0.844} & \textbf{0.829} \\
    BEV clustering (res \SI{0.30}{\metre})   & 0.993 & 0.640 & 0.779 & --- \\
    Higher MCS \num{20}                       & 0.976 & 0.755 & 0.852 & 0.801 \\
    Fragment merge (\SI{1.5}{\metre}, \num{30}\,pt) & 0.931 & 0.755 & 0.834 & 0.815 \\
    Adaptive HDBSCAN (\num{30}/\num{15}/\num{10}) & 0.935 & 0.748 & 0.831 & --- \\
    \bottomrule
  \end{tabular}
\end{table}

Each strategy fails for a structural reason, and the reasons are consistent with the
diagnosis of Section~\ref{sec:split-analysis}:

\begin{itemize}
  \item \textbf{Higher minimum cluster size (MCS \num{20})} is the most instructive.
    It is the only variant to \emph{improve} $F_1$ on sequence~00
    (\num{0.844}~$\rightarrow$~\num{0.852}), by discouraging the small splinter
    clusters. But it does not generalise: on the held-out sequence~08 it drops to
    $F_1$~\num{0.801}, \emph{below} the \num{0.829} baseline, because a larger
    minimum also erases genuinely sparse distant cars. The sequence-00 gain was
    over-fitting.
  \item \textbf{Post-clustering fragment merge} (join clusters within
    \SI{1.5}{\metre}) raises sequence-00 recall to \num{0.755} but loses precision
    (\num{0.984}~$\rightarrow$~\num{0.931}), because the same proximity rule also
    fuses nearby non-car structure --- walls, poles --- into oversized clusters. On
    held-out data the precision loss cancels the recall gain
    ($F_1$~\num{0.815}~$<$~\num{0.829}).
  \item \textbf{Distance-adaptive HDBSCAN} (a larger minimum cluster size near the
    ego, smaller far away) introduces boundary artefacts at the ring transitions and
    is worse than the global setting.
  \item \textbf{Bird's-eye-view clustering}, a cheaper
    connected-component alternative, collapses recall to \num{0.640}: projecting to
    two dimensions merges objects that overlap in the ground plane but are separated
    in height, and morphological cleanup destroys sparse distant clusters.
  \item \textbf{Temporal aggregation} --- accumulating object points over three
    consecutive frames before clustering, to densify sparse cars --- collapses the
    pipeline ($F_1$~\num{0.073}, recall \num{0.039}).\footnote{Source:
    \texttt{docs/session\_history.md} (2026-06 exploratory session). The
    \texttt{temporal\_window} feature was implemented and then fully reverted, so
    this result is on the pre-promotion baseline and was not re-run under the
    research freeze; it is reported here in prose, not in Table~\ref{tab:strategies},
    for that reason, and its value is the mechanism below rather than the exact
    figure.} Accumulating points does not densify
    the existing clusters; it reshapes the clustering landscape entirely
    (about \num{134000} accumulated points yield \num{645}+ clusters versus
    $\sim$\num{240} single-frame, of which fewer pass the geometric filters,
    \num{17} versus \num{34}).
  \item \textbf{Lowering the cluster-size thresholds} to admit sparser groups
    produced no change in true positives at all: the binding constraint is cluster
    quality, not the count threshold, so relaxing the threshold merely admits more
    noise (\texttt{docs/session\_history.md}).
\end{itemize}

The failure modes are strategy-specific --- over-fitting, precision loss, ring
artefacts, projection ambiguity, a reshaped clustering landscape --- but they share
a limit. The two strategies that reassemble after the split (fragment merge, and to
a lesser extent the proximity-based part of the adaptive scheme) run into the fact
that density-based clustering has no notion of objectness: once a dense car is
partitioned along an internal density gap, no proximity- or size-based rule can
reliably tell ``rejoin these two fragments of one car'' apart from ``do not join
these two adjacent cars.'' The strategies that instead change how clusters
\emph{form} --- a larger minimum size, a different clusterer, per-ring parameters,
pre-clustering accumulation --- avoid the reassembly trap but hit the same wall from
the other side: any single global density criterion that is loose enough to keep a
split car whole is also loose enough to fuse a close pair, so a sequence-00 gain
becomes a held-out loss. The one lever that helped, changing the clustering
\emph{resolution} (Section~\ref{sec:resolution-fix}), is also a formation change; it
works not by escaping this limit but by happening to sit at a resolution that closes
intra-car gaps before it starts closing inter-car ones. This is why the earlier
project conclusion --- that recall was fixed at a hard ceiling of about \num{0.74}
--- was reasonable given only the evidence
in Table~\ref{tab:strategies}, and also, as the next section shows, wrong: it read a
limit on \emph{these} repairs as a limit on the approach.

\subsection{Clustering resolution and the corrected ceiling}
\label{sec:resolution-fix}

The intervention that helped was not a repair at all but a change to the clustering
\emph{resolution}: re-voxelising the object cloud to a \SI{0.10}{\metre} grid
\emph{before} running HDBSCAN, rather than clustering the full-resolution
\SI{0.05}{\metre} cloud (adopted in the runtime optimisation of
Section~\ref{sec:results-runtime}). A coarser grid closes the small intra-car density gaps that HDBSCAN was
cutting along, so fewer splits form in the first place. It is a formation change like
the failed ones in Section~\ref{sec:repairs}, but it lands on the favourable side of
the trade described there --- coarse enough to close intra-car gaps, still fine
enough not to fuse neighbouring cars.

\begin{table}[H]
  \centering
  \caption[Effect of coarsening the clustering voxel]{Effect of coarsening the clustering voxel from \SI{0.05}{\metre} to
    \SI{0.10}{\metre} on the split rate (sequence~08, \num{100}
    frames, \num{810} ground-truth cars). More cars cluster cleanly and none newly
    merge --- the coarsening removes splits without fusing neighbours.}
  \label{tab:splitrate}
  \begin{tabular}{lcc}
    \toprule
    & Full-res (\SI{0.05}{\metre}) & Coarse (\SI{0.10}{\metre}) \\
    \midrule
    clustered ``ok''            & 504 (62.2\%) & 566 (69.9\%) \\
    split                        & 305 (37.7\%) & 242 (29.9\%) \\
    merged                       & 0 (0.0\%)    & 0 (0.0\%) \\
    mean clusters per split car  & 3.3          & 2.7 \\
    \bottomrule
  \end{tabular}
\end{table}

Table~\ref{tab:splitrate} shows the split rate falling by \num{7.8} points
(\SI{37.7}{\percent}~$\rightarrow$~\SI{29.9}{\percent}) with \emph{zero} new
merges.\footnote{The full-resolution counts here (\num{810} cars, split
\SI{37.7}{\percent}) and those in Table~\ref{tab:split} (\num{811}
cars, split \SI{37.1}{\percent}) are separate runs of the same analysis on the same
pre-promotion configuration; they differ by at most one car and \num{0.6} percentage
points, and the percentages are unaffected.}
At the pipeline level, on the full sequence~08, this lifts recall from \num{0.699} to
\num{0.730} ($+\num{1655}$ true positives) with precision essentially unchanged
(\num{0.903}~$\rightarrow$~\num{0.905}) and $F_1$ from \num{0.788} to \num{0.808}. The same coarsening also reduces clustering runtime, which is why it
was adopted; its detection benefit is the subject here and its runtime co-benefit is
reported in Section~\ref{sec:results-runtime}.

This result forces a correction to the earlier ceiling claim, which we state plainly
rather than bury. The ``$\sim$\num{0.74}
hard limit'' was \emph{partly a resolution artefact}, not a fundamental property of
the approach: the strategies of Section~\ref{sec:repairs} genuinely do fail, but none
of them tried this particular resolution setting, and changing it recovered part of
the loss. The standing recommendation is therefore
\texttt{cluster\_voxel\_size}~$=\SI{0.10}{\metre}$ in production, not ``accept
\num{0.74}.''

A smaller structural limit does remain, and it is architectural. Coarsening the voxel
further would eventually start merging adjacent cars, trading the split problem for a
merge problem; and the clustering-method benchmark of Section~\ref{sec:results-ablations}
shows that no fixed-radius density clusterer, at any tested radius, reaches even
HDBSCAN's recall. The ceiling that remains is the one imposed by density-based
clustering having no model of what a car is: it groups points by proximity and
density alone, so it cannot be pushed to the recall of a learned object detector by
tuning. Reaching that level would require replacing the clustering stage with a
learned detector, which is outside the scope of this thesis (Chapter~\ref{ch:conclusion}). The
value of this analysis is to have located the ceiling precisely --- it is in
the clustering stage, it is splitting rather than classification or merging, and it
is lower than first claimed --- and to have mapped which repairs cannot move it and
which one can.

\subsection{Recall analysis: summary}
\label{sec:recall-summary}

The recall shortfall is clustering-bound, not classifier-bound: the classifier and
track filter sit on top of a recall ceiling that geometric clustering already imposes
(Section~\ref{sec:recall-diagnosis}). The dominant mechanism is HDBSCAN splitting
large, close cars into fragments --- \SIrange{31}{37}{\percent} of cars, worst at
short range --- while merging is negligible, and the cleanly-clustered fraction sits
just below the achieved recall, so most split cars are lost
(Section~\ref{sec:split-analysis}). A search of six repair strategies --- higher
minimum cluster size, fragment merging, adaptive and BEV clustering, temporal
accumulation, and threshold lowering --- all fail to generalise, for a shared reason:
density-based clustering has no notion of objectness, so any criterion loose enough to
keep a split car whole also fuses adjacent cars, and reassembling fragments after the
split faces the same ambiguity (Section~\ref{sec:repairs}). The one effective lever was clustering
resolution: a coarser pre-clustering voxel prevents splits from forming, lifting
recall \num{0.699}~$\rightarrow$~\num{0.730} and correcting the earlier overstated
``hard ceiling'' (Section~\ref{sec:resolution-fix}). The residual limit is
architectural --- density clustering has no notion of objectness --- and is carried
to the limitations and future-work discussion of Chapter~\ref{ch:conclusion}.


%% ===== merged from sec_7_2_donor_metric.tex ===== %%

% Chamfer invalidity and the donor-metric validation battery (folds under Chapter 4).
%
% Scope: motivates why plain Chamfer is invalid on real cars (Finding #26), then
% defines and VALIDATES the donor metric. The completion coverage
% results it produces are reported in the completion-results section below. Honesty rules (personal_agenda
% P1): per-car medians / paired; static cars by construction; #37 presented as a
% metric-design lesson, not an erratum (claim C7). Numbers traceable to
% docs/completion/donor_metric.md and Findings #26/#32/#37.
%
With the detection front-end characterised, we now put the donor metric through its
validation battery and use it to carry the completion evaluation. The donor-frame metric it uses
to score completion was defined in Section~\ref{sec:donor-definition}. Because that
metric carries the completion claims, we first show it behaves --- passing an
explicit validation battery --- and then report what it, together with the
amodal-box evaluation, reveals about the completion output. The metric is defined
on static cars; moving cars are treated separately in
Section~\ref{sec:results-movers}.

\section{Validation of the donor metric}
\label{sec:donor-validation}

Before using it we required the metric to pass four pre-stated checks
(Table~\ref{tab:donor-gate}); it passes all four.

\begin{table}[H]
  \centering
  \caption[The four donor-metric validation checks]{The four validation checks for the donor metric, and the observed
    outcome (seq~08, per-car medians, $n=39$; the metric's original acceptance
    run, on the production path as it stood \emph{before} the
    Section~\ref{ch:completion-method} length prior). Coverage figures are the car-median
    \texttt{cov@0.1}; the guard figures are car-median out-of-box fractions.}
  \label{tab:donor-gate}
  \begin{tabular}{p{0.44\linewidth} p{0.48\linewidth}}
    \toprule
    Check & Outcome \\
    \midrule
    (a) The raw partial must rank \emph{last} on coverage (it covers no novel
      surface by construction). & Passes at every $\tau$: raw \num{0.000}
      $<$ mirrored \num{0.043} $<$ completed \num{0.304}. \\
    (b) Per-car scores must be stable enough for a per-car median to mean
      something. & Median per-car coverage IQR \num{0.14} --- moderate, as
      expected from varying view geometry; acceptable for medians. \\
    (c) The ranking must not depend on the threshold $\tau$. & Stable over
      $\tau \in \{0.10, 0.15, 0.20\}$: completed coverage
      \num{0.345} / \num{0.304} / \num{0.281}; the order never inverts. \\
    (d) Completion must not win by hallucinating: its out-of-box fraction must
      not exceed the mirrored baseline. & Passes: completed out-of-box
      \num{0.0003} $\leq$ mirrored \num{0.0083} (center/heading error makes the
      mirror spill more). \\
    \bottomrule
  \end{tabular}
\end{table}

Check~(d) is the one that makes the coverage number meaningful: without it, a
model could raise coverage simply by emitting surface in every direction. The
mirrored baseline in checks~(a) and~(d) is deliberately adversarial --- it is a
real way to add surface with zero learned parameters, so completion earns its
place only by beating it.

\subsection{The per-band hallucination guard}
\label{sec:donor-guard-lesson}

The guard as first designed --- check~(d) above --- was a single median pooled
over all cars, and that pooling hid a real failure mode. When we later evaluated
a geometry change (the fixed length prior of Section~\ref{ch:completion-method}) with this guard ---
now re-run on the promoted configuration, whose well-observed static set is
$n=40$ ($8+27+5$ across the compact/normal/long bands) --- the
pooled out-of-box fraction stayed comfortably low (\num{0.0020}) and the change
appeared safe. Splitting the same guard by car length told a different story
(Table~\ref{tab:d2}): on the \emph{compact} band the length prior over-extended
cars to \num{0.0433} out-of-box --- about \num{100}$\times$ that band's own
mirrored baseline --- while \num{32} of \num{40} cars sitting near \num{0.0015}
kept the pooled median quiet.

\begin{table}[H]
  \centering
  \caption[Pooled guard blind to a band-localised violation]{The pooled guard is blind to a band-localised violation. Out-of-box
    fraction as completed / that band's mirrored baseline, split by amodal-GT
    length. The final length prior passes the pooled check yet hallucinates on
    compacts at \num{100}$\times$ the band baseline.}
  \label{tab:d2}
  \begin{tabular}{lcccc}
    \toprule
    Config & compact $<$\SI{3.6}{\meter} & normal & long $\geq$\SI{4.6}{\meter} & pooled \\
    \midrule
    prior off        & 0.0003 / 0.0004 & 0.0008 / 0.0129 & 0.0000 / 0.0292 & 0.0004 \checkmark \\
    length prior on  & \textbf{0.0433 / 0.0004} & 0.0017 / 0.0129 & 0.0000 / 0.0292 & 0.0020 \checkmark \\
    \bottomrule
  \end{tabular}
\end{table}

Two lessons follow, and both are now built into the metric. First, the guard is
computed \emph{per length band against that band's own mirrored baseline}, not
as a pooled median; a change is only accepted if it clears the guard in every
band. Second, on the compact band the mirrored baseline is itself near zero
(\num{0.0004}): reflecting an already-short car barely leaves its box, so the
bar there is effectively ``do not leave the box at all,'' which \emph{no}
configuration clears. We therefore report the guard as a \emph{ratio} to the
band baseline rather than a pass/fail bit, because the ratio (for example, a
\num{16}$\times$ band violation reduced from an incumbent \num{100}$\times$ ---
a roughly \num{7}$\times$ reduction in hallucination) is the informative
quantity, while the binary bit is unclearable by construction and would hide
genuine progress.

This is a property of the measurement, not a retraction of a result: the
direction of the geometry change it was auditing still holds; what changed is
that the guard now reads at the resolution where the failure lives. Any future
change to the completion geometry is evaluated per band under this refined
guard.


%% ===== merged from sec_7_results.tex ===== %%

% Completion Evaluation results (folds under Chapter 4, after the donor-metric
% validation battery above).
%
% Scope: the completion RESULTS -- surface coverage (donor metric, #32); downstream
% box utility + length geometry (#29/#35/#36, #45 cross-ref); moving cars (#44);
% pre-registered held-out replication on seq 00 (#42); summary + limitations. The
% metric itself is defined in the validation-battery section above; the method in
% Chapter 3 (Methodology).
%
% Honesty rules (personal_agenda P1): per-car medians / paired Wilcoxon; static
% cars by construction; trade-offs stated plainly (per-car length estimate costs
% donor coverage); negatives (#45) as findings, not apologies; measured/inferred
% marked. Numbers traceable to docs/findings.md #29/#32/#35/#36/#37/#41/#42/#44/#45,
% docs/completion/donor_metric.md, and docs/plans/t9{b,c}_*heldout_seq00.md.
%
% Figures (via \graphicspath into the gitignored output/ tree; regenerable):
%   Fig 7.1 donor_metric_08.png              (donor_metric_viz.py)
%   Fig 7.2 completion_box_overlays_08.png   (completion_box_eval_viz.py)
%   Fig 7.3 t11_mover_completions_bev.png    (t11_mover_plausibility.py)
% Floats set [htbp].
%
Section~\ref{sec:donor-metric} established a metric that measures, on real cars, how much
previously-unseen surface a completion recovers, and showed that it passes an
explicit validation battery. This part of the chapter uses it. It reports three
things the completion method was built to deliver and one it was not: that the
completion adds genuinely unseen surface (Section~\ref{sec:results-coverage}),
that this translates into a better amodal bounding box (Section~\ref{sec:results-box}),
that moving cars --- which the static-car metrics cannot score --- complete no
worse than static ones (Section~\ref{sec:results-movers}), and that both headline
claims survive a pre-registered replication on a held-out sequence
(Section~\ref{sec:results-heldout}). Section~\ref{sec:results-summary} draws the
limitations together. Throughout, the unit of analysis is the car: because the
frames of one parked car are autocorrelated, each car is reduced to its per-car
median and significance is tested across cars with the Wilcoxon signed-rank test,
the paired design of Section~\ref{sec:donor-metric}.

\section{Real-data surface coverage}
\label{sec:results-coverage}

The primary completion result is that PCN reconstructs surface the input frame
never observed. Table~\ref{tab:donor-coverage} reports the donor metric on \num{39} of seq~08's
\num{40} well-observed static cars (one car's pairs are all gate-skipped), scoring
three clouds --- the raw partial, its symmetry mirror, and the completion ---
against each car's
donor-only novel set. This is the metric's acceptance run, taken
on the production path \emph{before} the Section~\ref{ch:completion-method} length prior; the length
geometry is evaluated in Section~\ref{sec:results-box}.

\begin{table}[H]
  \centering
  \caption[Real-data surface coverage on the donor metric]{Real-data surface coverage on the donor metric (seq~08, per-car
    medians, $n=39$ of the \num{40} well-observed static cars --- one car's pairs
    are all gate-skipped; $\tau=\SI{0.15}{\meter}$). \texttt{cov@0.1} is the fraction of never-observed (donor-only)
    surface recovered within \SI{0.1}{\meter}; median novel distance is its
    threshold-free companion; out-of-box is the hallucination guard. All three
    pairwise contrasts are significant at $p<10^{-6}$.}
  \label{tab:donor-coverage}
  \begin{tabular}{lccc}
    \toprule
    Cloud & \texttt{cov@0.1} $\uparrow$ & med.\ novel dist.\ (m) $\downarrow$ & out-of-box $\downarrow$ \\
    \midrule
    raw partial       & 0.000          & 0.518          & 0.000 \\
    mirrored          & 0.043          & 0.332          & 0.008 \\
    \textbf{completed} & \textbf{0.304} & \textbf{0.161} & \textbf{0.000} \\
    \bottomrule
  \end{tabular}
\end{table}

The completion recovers about \SI{30}{\percent} of never-observed surface within
\SI{0.1}{\meter} --- \emph{seven times} the symmetry-mirror baseline --- while its
out-of-box fraction stays at the raw partial's floor, so the coverage is real
surface, not surface invented everywhere. The mirror is the informative
comparison: reflecting the input across the car's symmetry plane is a real,
zero-parameter way to add surface, and the learned model earns its place only by
beating it, which it does more than sixfold. This is the project's first positive
real-data evidence that completion adds unseen structure, and it complements the
box-level value of Section~\ref{sec:results-box}.

Two breakdowns qualify the headline. First, coverage is uneven over the car: the occluded \emph{far side} is well recovered (completed cov
\num{0.321}) and the roof reasonably so (\num{0.203}), but the \emph{far end} ---
the portion of the car truncated by the single viewpoint --- is the weakest
region (\num{0.133}). The far end is exactly where a one-sided partial is
normalised around its observed near portion, and it is what the length geometry
of Section~\ref{sec:results-box} sets out to move. (The mirror scores essentially
zero on the far end by construction --- a reflection across the length plane
cannot add end surface --- which is a built-in sanity check on the region split.)
Second, the coverage is robust to input sparsity: completed \texttt{cov@0.1}
holds at \num{0.30}--\num{0.33} from the sparsest (${<}100$ points) to the densest
(${\geq}300$) inputs, while the mirror falls from \num{0.083} on the densest
inputs to \num{0.010} on the sparsest, as the reflection of a sparse cluster
leaves ever more of the car uncovered.
Completion helps most, relative to the cheap baseline, exactly where the sensor
saw least.

\begin{figure}[htbp]
  \centering
  \includegraphics[width=\linewidth]{donor_metric_08.png}
  \caption[The donor metric on six seq-08 cars]{The donor metric in bird's-eye view on six seq-08 cars, spanning the
    coverage range from best (top-left) to a residual failure (bottom-right). The
    \emph{novel set} --- the donor-only surface the input frame never saw --- is
    drawn blue where the completion recovers it within \SI{0.1}{\meter} and red
    where it does not, so the blue fraction of the target is exactly the reported
    \texttt{cov@0.1} in each title. The three candidate clouds are shown faintly
    for context: the input partial (grey), its symmetry mirror (orange), and the
    completion itself (green). The upper panels close the occluded far side almost
    entirely; the last isolates the residual failure modes --- far-end
    under-completion and heading/centre error on diagonal or sparse views --- that
    Section~\ref{sec:results-box} and Section~\ref{ch:completion-method} address.
    Rendered by \texttt{scratchpad/donor\_metric\_viz.py}.}
  \label{fig:donor-coverage}
\end{figure}

\section{Downstream utility: the amodal bounding box}
\label{sec:results-box}

Surface coverage shows the completion is geometrically faithful; the box
evaluation shows it is \emph{useful}. For each true-positive (frame, car) pair on
the well-observed static cars we fit an oriented box to the raw partial and to the
completion with the \emph{same} fitter used to build the amodal ground-truth
boxes (so fitter bias cancels), and compare each against the GT box. This is the
pre-registered ``completion adds value'' test.
Table~\ref{tab:box-quality} reports it.

\begin{table}[H]
  \centering
  \caption[Amodal box quality: raw partial vs.\ completion]{Amodal box quality, raw partial vs.\ completion (seq~08, per-car
    medians, $n=39$ of the \num{40} well-observed static cars --- one has no
    completed pair). BEV IoU is the oriented
    bird's-eye-view box overlap; $|\Delta W|,|\Delta H|,|\Delta L|$ are absolute
    dimension errors; centre error is in the ground plane. $p$ is the Wilcoxon
    signed-rank across cars. These figures use the original detection
    configuration; the width, height, and overlap wins reproduce under the promoted
    configuration of Table~\ref{tab:length-geometry} (BEV IoU
    \num{0.725}$\to$\num{0.747}, $|\Delta W|$ \num{0.210}$\to$\num{0.122},
    $|\Delta H|$ \num{0.227}$\to$\num{0.144}).}
  \label{tab:box-quality}
  \begin{tabular}{lcccl}
    \toprule
    Metric & Raw & Completed & Wilcoxon $p$ & \\
    \midrule
    BEV IoU $\uparrow$        & 0.707 & \textbf{0.747} & \num{1.9e-3}  & better \\
    $|\Delta W|$ (m) $\downarrow$ & 0.270 & \textbf{0.170} & \num{1.5e-4}  & better \\
    $|\Delta H|$ (m) $\downarrow$ & 0.255 & \textbf{0.131} & \num{1.6e-10} & better \\
    centre err.\ (m) $\downarrow$ & 0.286 & \textbf{0.234} & \num{2.8e-5}  & better \\
    $|\Delta L|$ (m) $\downarrow$ & 0.447 & 0.456          & 0.65          & neutral \\
    yaw err.\ (deg) $\downarrow$  & 3.5   & 3.0            & 0.74          & neutral \\
    \bottomrule
  \end{tabular}
\end{table}

Completion improves the box on width, height, centre, and overall BEV overlap
with strong significance, and nothing degrades significantly. The improvements
are largest where the sensor saw least --- BEV IoU by input size rises
\num{0.461}$\to$\num{0.599} on the sparsest cars (${<}100$ points),
\num{0.585}$\to$\num{0.678} at \num{100}--\num{300}, and \num{0.703}$\to$\num{0.744}
above \num{300} --- which is the right profile for the use case, since a
well-observed car needs little help. The width gain is not an artefact of the
GT accumulation: it holds both for the \num{13} cars whose width the accumulation
genuinely constrains (both flanks eventually seen: \num{0.291}$\to$\num{0.182})
and for the other \num{26} (\num{0.167}$\to$\num{0.118}).

\paragraph{Length was the exception, and the geometry work fixed it.} In
Table~\ref{tab:box-quality} length is neutral --- the completion neither helps nor
hurts $|\Delta L|$. That was too generous a reading of an untreated defect: the
completion of a partial truncated at the occluded far end stops short, so it
\emph{under}-completes length. Table~\ref{tab:length-geometry} isolates this on
the promoted configuration ($n=40$) and tracks the two Section~\ref{ch:completion-method} fixes ---
the fixed \SI{4.14}{\meter} prior and the per-car length estimate --- against the raw partial and the completion \emph{without} any
length correction.

\begin{table}[H]
  \centering
  \caption[Length geometry on the amodal box]{Length geometry on the amodal box (seq~08, well-observed static cars,
    per-car medians, $n=40$, promoted configuration).
    ``no prior'' is the completion with the canonical-frame centre but no
    longitudinal push; ``fixed \SI{4.14}{}'' is the population-median prior;
    ``per-car'' is the final \nth{90}-percentile track estimate. The compact
    row is signed $\Delta L$ on cars ${<}\SI{3.6}{\meter}$ ($n=8$), where a fixed
    prior over-extends.}
  \label{tab:length-geometry}
  \begin{tabular}{lcccc}
    \toprule
    Metric & Raw & no prior & fixed \SI{4.14}{} & \textbf{per-car (final)} \\
    \midrule
    $|\Delta L|$ (m) $\downarrow$   & 0.428 & 0.476 & 0.354 & \textbf{0.304} \\
    BEV IoU $\uparrow$              & 0.725 & 0.743 & 0.747 & \textbf{0.771} \\
    centre err.\ (m) $\downarrow$   & 0.271 & 0.227 & 0.229 & \textbf{0.184} \\
    \midrule
    compact signed $\Delta L$ (m)   & $-0.007$ & --- & $+0.295$ & \textbf{$+0.063$} \\
    \bottomrule
  \end{tabular}
\end{table}

Two things are visible here, and both are stated as the honest picture rather
than the flattering one. First, Table~\ref{tab:box-quality}'s neutral
$|\Delta L|$ averages over- and under-completion and so hides a real regression:
the \emph{signed} length shows the completion systematically \emph{under}-extends
the occluded far end (normal cars $-0.485\to-0.545$). On the
promoted configuration this surfaces even in the absolute $|\Delta L|$, where
completion \emph{without} a length prior makes box length \emph{worse} than the
raw partial (\num{0.428}$\to$\num{0.476}, $p=0.027$), which the far-end truncation
causes. The fixed prior
reverses it (\num{0.354}) and the per-car estimate improves it further
(\num{0.304}) while lifting BEV IoU to \num{0.771} and cutting centre error to
\num{0.184}. Second, moving from the fixed prior to the per-car estimate is what
removes the over-extension of genuinely compact cars: their signed $\Delta L$ goes
from \num{+0.295} (the fixed median over-extends a short car) to \num{+0.063}
($p=0.016$ at $n=8$), with the same fix showing up in the per-band hallucination
guard of Section~\ref{sec:donor-metric}.

\paragraph{The trade-off, stated plainly.} The per-car estimate buys this box
accuracy at a measured cost in donor coverage: pooled completed \texttt{cov@0.1}
falls \num{0.403}$\to$\num{0.364} and far-end coverage \num{0.346}$\to$\num{0.316}
relative to the fixed prior. Per band it is better on compact and
long cars and worse on normal ones. The two metrics genuinely disagree here, and
a control run settled which reading is right: deliberately over-extending
(\nth{90} percentile $+\SI{0.45}{\meter}$) improved \emph{both} coverage and the
box on normal and long cars, which means those completions are still honestly too
short rather than the metric being fooled. In other words the per-car estimate
trades a little pooled coverage --- driven by the normal band it slightly
shortens --- for a correctly-sized box across all bands, and the residual
shortfall is a real limitation, addressed next.

\paragraph{A bounded negative on the residual shortfall.} Targeting each car's
\emph{true} length still leaves completions short on normal and long cars, which
pointed to a second under-extension mechanism inside the completion itself (PCN
under-fills its normalised frame, and the centre push doubles as an output
rescale). Section~\ref{ch:completion-method} investigated this under pre-registration (the
``closed negative''): decoupling the scale from the length push and adding a
calibrated fill factor did fix the target (normal-car $|\Delta L|$
\num{0.329}$\to$\num{0.232}) but over-extended compacts and failed the
pre-registered compact non-regression guard on both test sequences. The verdict
was to reject it; the knobs stay in the code defaulted off. It is reported here
because it \emph{bounds} the result: the residual far-end shortfall on long cars
is a characterised, length-dependent limit, not an oversight.

\begin{figure}[htbp]
  \centering
  \includegraphics[width=\linewidth]{completion_box_overlays_08.png}
  \caption[Amodal box overlays on six seq-08 cars]{Amodal box overlays on six seq-08 cars: the ground-truth box (black),
    the box fitted to the raw single-frame partial (blue), and the box fitted to
    the completion (green), over the raw points (grey). The completion tightens
    the box on the occluded width and far end and centres it better; the one
    strong regression panel is a heading error on a sparse (\num{109}-point)
    input, a diagonal-view failure mode also visible in
    Figure~\ref{fig:donor-coverage}. Rendered by
    \texttt{scratchpad/completion\_box\_eval\_viz.py}.}
  \label{fig:box-overlays}
\end{figure}

\section{Moving cars}
\label{sec:results-movers}

Both metrics so far are defined on \emph{static} cars by construction: a paired
amodal reference needs a stationary target. Production nonetheless completes
moving cars, and those completions were never validated. Because completion runs
on a single reference frame's cluster --- not on accumulated points --- there is
no motion smear in its input, so there is a mechanistic reason to expect motion
not to matter; this section checks it. On the final seq-08 output
(\num{518} completed car tracks) each track is labelled static, moving, or
ambiguous by the net horizontal displacement of its centroid (a label-free proxy,
since the output tracks are unlabelled), and each completion is scored by a
size-plausibility recipe (bird's-eye dimensions falling in the car range).
Table~\ref{tab:movers} gives the result.

\begin{table}[H]
  \centering
  \caption[Completion plausibility by track motion]{Completion plausibility by track motion (final seq-08 output,
    \num{518} completed car tracks). Motion is the net displacement
    of the track centroid (median of first vs.\ last five frames). Plausible =
    bird's-eye dimensions in the car range ($L\in[3.3,4.9]$, $W\in[1.5,2.1]$,
    $H\in[1.1,1.7]$~m).}
  \label{tab:movers}
  \begin{tabular}{lcccc}
    \toprule
    Group (net displacement) & $n$ & plausible & rate & median $L/W/H$ (m) \\
    \midrule
    static ($\leq\SI{2}{\meter}$)   & 419 & 225 & \SI{53.7}{\percent} & 3.81 / 1.94 / 1.49 \\
    \textbf{moving ($\geq\SI{5}{\meter}$)} & \textbf{19} & \textbf{11} & \textbf{\SI{57.9}{\percent}} & 3.84 / 2.05 / 1.44 \\
    ambiguous (\SIrange{2}{5}{\meter}) & 80  & 46  & \SI{57.5}{\percent} & 3.95 / 1.92 / 1.49 \\
    all                             & 518 & 282 & \SI{54.4}{\percent} & 3.83 / 1.94 / 1.49 \\
    \bottomrule
  \end{tabular}
\end{table}

Moving cars complete \emph{no worse} than static ones --- slightly higher, in
fact (\SI{57.9}{} vs.\ \SI{53.7}{\percent}), well within noise at $n=19$ --- with
near-identical median dimensions, and every motion bucket clusters at
\SIrange{54}{58}{\percent}. Motion does not degrade completion plausibility, as
the single-frame input predicts. One honest caveat about the absolute level, not
the comparison: the plausibility recipe uses axis-aligned dimensions, which
overestimate width and length for cars oriented diagonally to the world axes, so
the \raisebox{0pt}[0pt][0pt]{$\sim$}\SI{54}{\percent} rate \emph{understates}
completion quality (most failures are inflated width on genuine car-shaped
footprints). Since both groups use the identical recipe, the static-vs-moving
conclusion is unaffected. The value of the check is that the statics-only design
of the earlier sections does not hide a mover failure mode.

\begin{figure}[htbp]
  \centering
  \includegraphics[width=\linewidth]{t11_mover_completions_bev.png}
  \caption[Completions of moving-car tracks on seq~08]{Bird's-eye completions of moving-car tracks on seq~08 (blue partial
    input, green completion). Most form recognisable car footprints; the failures
    are width inflation on diagonally-oriented cars, an artefact of the
    axis-aligned plausibility measure rather than of motion. Rendered by
    \texttt{scratchpad/t11\_mover\_plausibility.py}.}
  \label{fig:movers}
\end{figure}

\section{A pre-registered held-out replication}
\label{sec:results-heldout}

The completion constants --- the width and length priors, the scale correction,
the gate thresholds --- were all developed on seq~08. A fair test of whether the
method \emph{generalises} therefore has to run on a sequence that took no part in
that development, under a protocol fixed in advance so the outcome cannot be
reinterpreted after the fact. We pre-registered one (\texttt{preregistration\_heldout.md},
committed with an immutable timestamp before any held-out run). The held-out
sequence is seq~00 (seq~05 had too few well-observed
static cars to test); the two \emph{primary} metrics were fixed as BEV box IoU
(Section~\ref{sec:results-box}) and donor \texttt{cov@0.1}
(Section~\ref{sec:results-coverage}), each requiring the completion to beat the
raw partial with $p<0.05$; three refutation rules (fails to generalise / seq-08
specific length constants / an uncovered result) and a three-level outcome
taxonomy were registered alongside. Nothing was retuned for seq~00.

One confound is logged rather than hidden: seq~00 is in the classifier's training
split, so its \emph{detection} recall reads optimistically. This does not touch
the completion claim, which is a \emph{paired} raw-vs-completed comparison on
identical true-positive inputs, and PCN was trained only on synthetic data and
never saw seq~00. Table~\ref{tab:heldout} gives the outcome on the \num{45}
well-observed static cars.

\begin{table}[H]
  \centering
  \caption[Held-out replication on seq~00]{Held-out replication on seq~00 (well-observed static cars, per-car
    medians, $n=45$; the two pre-registered primaries in bold, secondary
    dimensions below). Both primaries are decisive wins over the
    raw partial.}
  \label{tab:heldout}
  \begin{tabular}{lccl}
    \toprule
    Metric & Raw & Completed & Wilcoxon $p$ \\
    \midrule
    \textbf{BEV IoU} $\uparrow$          & \textbf{0.739} & \textbf{0.766} & \textbf{\num{1.6e-3}} \\
    \textbf{donor \texttt{cov@0.1}} $\uparrow$ & \textbf{0.000} & \textbf{0.413} & \textbf{${\approx}0$} \\
    $|\Delta W|$ (m) $\downarrow$        & 0.203 & 0.165 & \num{4.0e-4} \\
    $|\Delta H|$ (m) $\downarrow$        & 0.278 & 0.120 & \num{7.8e-9} \\
    \bottomrule
  \end{tabular}
\end{table}

Both headline claims replicate: on a sequence that played no part in tuning, the
completion beats the raw partial on box overlap (\num{0.739}$\to$\num{0.766},
$p=\num{1.6e-3}$) and on surface coverage (\num{0.000}$\to$\num{0.413},
$p\approx0$), with the secondary width and height errors improving as they did on
seq~08. Applying the pre-registered rules verbatim, the arbitration verdict was
\textbf{partially holds}: neither generalisation rule triggered ---
both primaries won and no testable length band regressed against its seq-08 level
--- but seq~00's well-observed set contains \emph{no} long cars
(${\geq}\SI{4.6}{\meter}$), so the long-band predictions and the long-band
hallucination guard are simply untestable on it. The downgrade from ``holds'' to
``partially holds'' is therefore a \emph{coverage gap in the held-out data, not a
weak result}: everything that could be tested passed, and the one thing that could
not be tested is a property of seq~00's car population, not of the method.

\section{Summary and limitations}
\label{sec:results-summary}

The completion method rests on three measured results. It adds genuinely unseen
surface --- about \SI{30}{\percent} of never-observed surface within
\SI{0.1}{\meter}, seven times the symmetry-mirror baseline, with essentially no
hallucination (Section~\ref{sec:results-coverage}). It improves the amodal box on
width, height, centre, and overall overlap, most where the sensor saw least, and
the length-geometry work turns an untreated length regression into a net gain
(Section~\ref{sec:results-box}). And both headline claims survive a
pre-registered replication on a held-out sequence
(Section~\ref{sec:results-heldout}). Moving cars, which the static-car metrics
cannot score, complete no worse than static ones
(Section~\ref{sec:results-movers}).

The limitations are stated as findings, each already quantified. The
\emph{far end} of the car is the weakest-covered region, and on long cars the
residual under-extension is a characterised, length-dependent limit that a
pre-registered fix could not remove without regressing compacts.
The per-car length estimate falls back to the fixed population prior on
\SI{23}{\percent} of final tracks --- those with fewer than five gate-passed
frames --- which is above a comfortable threshold and is not
solved. The sparsest inputs complete more coarsely, and the axis-aligned plausibility
measure inflates the \emph{reported} width of diagonally-oriented cars. And the two quantitative metrics measure static cars directly;
moving cars are covered only by the weaker plausibility check. None of these
overturns the headline --- completion adds unseen surface and a better box, and
generalises --- but they mark exactly where the method is soft, and they set the
agenda for the discussion that follows.
